\chapter{Marco operativo y normativo del dominio de emergencias}
\label{cap:3}

En este capítulo se analiza en detalle el entorno regulatorio y operativo que condiciona el diseño del prototipo de priorización. En primer lugar, se expone el marco legal de la protección civil a nivel estatal y autonómico, con especial énfasis en la legislación de Castilla y León. Posteriormente, se describen las variables operativas y las reglas que sustentan la asignación académica de prioridades, justificando su alineación con las normas vigentes. Por último, se examinan los aspectos regulatorios relacionados con la privacidad, la protección de datos personales y la necesidad ineludible de supervisión humana en la toma de decisiones críticas.

\section{Finalidad del capítulo}

El sistema de apoyo a la decisión propuesto en este Trabajo Fin de Máster (TFM) se sitúa en un dominio especialmente regulado: la atención de emergencias, la protección civil y la gestión de información sensible en los primeros minutos de un incidente. Por ello, antes de describir la arquitectura técnica del prototipo resulta necesario fijar el marco operativo y normativo que justifica sus variables, sus reglas de prioridad y sus límites de uso. El cumplimiento de estos principios legales no solo responde a imperativos éticos y de gobernanza, sino que también determina de forma directa la viabilidad de la integración del sistema en centros de coordinación real.

Este capítulo no pretende realizar un estudio jurídico exhaustivo. Su función es más concreta: extraer de la normativa estatal, autonómica y europea los criterios que pueden convertirse en decisiones de diseño. En particular, interesa identificar qué elementos normativos respaldan la detección de riesgo vital, la valoración de víctimas, la vulnerabilidad, la progresión del incidente, la coordinación entre recursos, la trazabilidad de las recomendaciones, la privacidad y la supervisión humana. De este modo, se tiende un puente metodológico sólido entre los principios de la jurisprudencia de emergencias y la ingeniería de características del prototipo.

La línea de trabajo adoptada tras el rediseño del proyecto se centra en Castilla y León. El caso empírico se apoya en el Registro de Emergencias del 112 de Castilla y León, correspondiente al periodo 2008--2022, y el marco autonómico principal queda definido por la Ley 4/2007, de Protección Ciudadana de Castilla y León, el Plan Territorial de Protección Civil de Castilla y León (PLANCAL) y los planes especiales incorporados al corpus normativo del prototipo \cite{Ley4_2007_CYL,PLANCAL_DEC_4_2019,Registro112CyL}. Sobre esta base se construye una escala académica P1--P4, no oficial, orientada a priorizar incidentes de forma explicable y auditable.

\section{Sistema español de protección civil}

La protección civil en España se configura como un sistema público orientado a proteger a las personas, los bienes y el medio ambiente ante situaciones de grave riesgo, catástrofe o calamidad pública \cite{Ley17_2015}. Su organización combina competencias estatales, autonómicas y locales, de modo que la respuesta a una emergencia no depende de una única autoridad ni de un único centro operativo, sino de una red de estructuras de mando, coordinación y apoyo técnico.

Desde la perspectiva del sistema propuesto, esta configuración tiene tres consecuencias. En primer lugar, la prioridad de un incidente no puede deducirse solo de su categoría nominal: un accidente de tráfico, un incendio o una incidencia sanitaria pueden exigir tratamientos muy distintos según víctimas, localización, propagación, vulnerabilidad o simultaneidad con otros avisos. En segundo lugar, la respuesta se produce bajo incertidumbre: durante los primeros minutos el operador dispone de texto libre, categoría preliminar, localización aproximada y señales parciales, no de una fotografía completa de la emergencia. En tercer lugar, cualquier recomendación automatizada debe ser subordinada a la decisión humana, porque el sistema no moviliza recursos ni declara situaciones operativas.

La Ley 17/2015 aporta el marco conceptual general: riesgo, emergencia, afectación a personas, protección de bienes, servicios esenciales, deberes de colaboración y coordinación entre administraciones \cite{Ley17_2015}. Estos conceptos justifican que el prototipo no se limite a detectar palabras graves en el texto, sino que estructure la información en variables operativas: riesgo vital, número estimado de víctimas, gravedad de lesiones, población vulnerable, emplazamiento crítico, multirriesgo o accesibilidad de recursos.

La Norma Básica de Protección Civil, aprobada mediante el Real Decreto 524/2023, refuerza esta lógica al ordenar la planificación por riesgos, establecer criterios comunes para los planes y situar la gestión de emergencias dentro de una arquitectura graduada de respuesta \cite{RD524_2023}. Para este TFM, su utilidad no consiste en trasladar mecánicamente las situaciones operativas al modelo, sino en justificar que la prioridad debe entenderse como una gradación: hay incidentes leves, incidentes relevantes, incidentes graves e incidentes críticos, y el sistema debe distinguirlos con criterios reproducibles.

El Plan Estatal General de Emergencias de Protección Civil (PLEGEM) completa el marco estatal al describir la coordinación ante emergencias de interés nacional o de dimensión supraautonómica \cite{PLEGEM}. Aunque el prototipo no pretende activar mecanismos estatales, el PLEGEM resulta útil para modelar señales de complejidad: multirriesgo, avisos simultáneos, insuficiencia potencial de medios ordinarios o necesidad de coordinación ampliada.

\section{Marco autonómico de Castilla y León}

Castilla y León constituye el ámbito territorial de referencia del prototipo por dos razones complementarias. La primera es empírica: existe un registro público de emergencias del 112 que permite trabajar con incidentes reales y evitar una validación basada exclusivamente en ejemplos simulados \cite{Registro112CyL}. La segunda es normativa: la comunidad autónoma dispone de una ley propia de protección ciudadana y de planes territoriales y especiales que permiten anclar la interpretación operativa de las variables del sistema.

\subsection{Ley 4/2007 de Protección Ciudadana de Castilla y León}

La Ley 4/2007, de Protección Ciudadana de Castilla y León, proporciona el marco autonómico general para la organización de la protección civil, la coordinación de emergencias y la actuación de los servicios públicos en la comunidad \cite{Ley4_2007_CYL}. Para el sistema propuesto, su interés principal es que sitúa la atención de emergencias dentro de un esquema institucional donde la rapidez de respuesta, la coordinación y la información al ciudadano son elementos centrales.

Desde el punto de vista del modelado, esta ley respalda variables relacionadas con la inmediatez y la capacidad operativa, especialmente la accesibilidad de recursos (V15), la existencia de emplazamientos críticos (V07) y la necesidad de registrar la actuación para auditoría posterior. También refuerza una idea metodológica importante: el sistema recomienda prioridades, pero la decisión final corresponde siempre a personal humano integrado en la organización competente.

\subsection{PLANCAL como ancla de la escala P1--P4}

El Decreto 4/2019 aprueba el Plan Territorial de Protección Civil de Castilla y León (PLANCAL), pieza central para conectar el marco normativo con la escala académica de prioridad utilizada en este TFM \cite{PLANCAL_DEC_4_2019}. PLANCAL no define una etiqueta P1--P4 para el dataset, ni el proyecto afirma que exista una equivalencia oficial. Su aportación consiste en ofrecer una lógica de fases, situaciones y coordinación que permite construir una guía de etiquetado razonada.

La escala P1--P4 del prototipo debe entenderse como una abstracción funcional. P1 representa incidentes con riesgo vital, víctimas graves, atrapamiento, intoxicación peligrosa, propagación severa o necesidad de atención inmediata. P2 agrupa incidentes relevantes con potencial de deterioro o afectación significativa, pero sin la criticidad inmediata de P1. P3 recoge situaciones que requieren gestión ordinaria o seguimiento, y P4 corresponde a incidentes de baja urgencia operativa o sin afectados. Esta escala se utiliza para entrenamiento, evaluación y explicación, pero no sustituye categorías oficiales de mando.

\subsection{Planes especiales: INFOCAL, INUNCYL y MPCyL}

Los planes especiales de Castilla y León aportan criterios sectoriales que enriquecen la interpretación de incidentes concretos. INFOCAL, aprobado por el Decreto 6/2025, resulta relevante para incendios forestales, propagación, amenaza a núcleos habitados, simultaneidad de focos y condiciones que pueden elevar la prioridad de una incidencia inicialmente limitada \cite{INFOCAL_DEC_6_2025}. En el prototipo, este marco se relaciona con V12 (\texttt{riesgo\_propagacion}) y con señales textuales como \texttt{signal\_incendio}.

INUNCYL, plan de protección civil ante el riesgo de inundaciones en Castilla y León, aporta el marco conceptual para valorar episodios de inundación, avenidas, desbordamientos y afectación territorial \cite{INUNCYL}. En la versión v0.1.0 del proyecto, las variables de enriquecimiento meteorológico e hidrológico asociadas a INUNCYL se difieren a v0.2.0, pero la norma permanece en el marco teórico y en el corpus RAG porque justifica una línea natural de extensión.

MPCyL, plan de protección civil ante el riesgo de transportes de mercancías peligrosas, es relevante para incidentes con sustancias químicas, cisternas, fugas, intoxicaciones o riesgos tecnológicos en transporte \cite{MPCYL_ACUERDO_3_2008}. En v0.1.0 se utiliza principalmente como anclaje normativo de reglas asociadas a \texttt{signal\_intoxicacion} y a escenarios de mercancías peligrosas. La planificación reconoce una limitación práctica: el corpus dispone del decreto de aprobación del plan, pero no del documento técnico completo del MPCyL. Por tanto, esta fuente se documenta como limitación específica del RAG.

Para mitigar parcialmente esta carencia, el prototipo incorpora una inyección sintética controlada de anclajes normativos asociados a mercancías peligrosas. Esta técnica no sustituye al texto técnico completo ni se presenta como fuente jurídica autónoma; su función es mantener trazabilidad mínima entre las reglas de la Capa 2, los riesgos químicos o de transporte y el decreto de aprobación disponible. La incorporación del plan técnico completo queda como tarea prioritaria para una versión posterior del corpus normativo.

\section{Normativa de datos, comunicaciones e inteligencia artificial}

El dominio no solo exige priorizar bien; exige hacerlo de forma trazable, privada y compatible con la supervisión humana. Por este motivo, el marco normativo del proyecto incluye también normas sobre comunicaciones de emergencia, protección de datos e inteligencia artificial.

La Ley 11/2022, General de Telecomunicaciones, resulta relevante por el papel de los servicios de emergencia, la localización de llamadas y el tratamiento de información asociada a comunicaciones \cite{Ley11_2022}. En el prototipo, esta norma no se traduce en una integración técnica con sistemas de telecomunicaciones, sino en una cautela de diseño: la localización se trata como dato operativo sensible y solo se conserva o transforma cuando aporta valor a la priorización y a la trazabilidad.

El Reglamento General de Protección de Datos (RGPD) y la Ley Orgánica 3/2018 (LOPDGDD) imponen principios de minimización, limitación de finalidad, seguridad y responsabilidad proactiva \cite{RGPD,LOPDGDD}. Estos principios justifican el rechazo de variables innecesarias, el uso de hashes en los logs de inferencia, la separación entre datos de entrada y resultados del sistema, y la exclusión de campos que puedan revelar decisiones posteriores o resultados clínicos. Esta restricción conecta directamente con la política anti-\textit{leakage}: el sistema no debe aprender a priorizar usando información que en la realidad solo estaría disponible después de la decisión inicial.

El Reglamento (UE) 2024/1689 de Inteligencia Artificial (UE AI Act) introduce un marco de especial relevancia para sistemas aplicados a servicios esenciales, emergencias y apoyo a decisiones con posible impacto sobre personas \cite{EUAIAct2024}. El TFM no afirma que el prototipo sea un producto certificado ni listo para despliegue. Su posición es más prudente: el diseño se alinea con requisitos de alto nivel como supervisión humana, trazabilidad, transparencia, gestión del riesgo, calidad de datos y robustez. Por ello, la arquitectura se formula como sistema de apoyo a la decisión: recomienda, explica y registra, pero no decide ni moviliza recursos de forma autónoma.

\section{Situaciones operativas y escala académica P1--P4}

Las situaciones operativas de protección civil ordenan el mando, la coordinación y la movilización de recursos ante emergencias. Su finalidad es institucional y organizativa. La prioridad P1--P4 del sistema, en cambio, es una escala académica diseñada para ordenar incidentes por urgencia relativa en el contexto del prototipo.

Esta diferencia debe quedar clara para evitar una confusión metodológica. Una situación operativa puede referirse al estado global de una emergencia en un territorio, mientras que P1--P4 se asigna a un incidente concreto. Un accidente con una persona atrapada puede ser P1 aunque no exista una situación territorial de emergencia elevada. Del mismo modo, un escenario de prealerta territorial puede contener avisos individuales de baja prioridad. La relación entre ambas escalas es, por tanto, inspiracional y funcional, no jurídica ni uno a uno.

La escala P1--P4 se construye con tres objetivos. Primero, permitir que los tres autores generen etiquetas débiles comparables sobre el dataset real. Segundo, entrenar y evaluar un motor interpretable, basado en RuleFit y comparado con un baseline experto. Tercero, generar explicaciones comprensibles para el operador, citando reglas activadas y normas de referencia. Esta formulación mantiene el sistema dentro de la línea de trabajo del proyecto: núcleo explicable, supervisión humana y trazabilidad normativa.

\section{Taxonomía operativa de incidentes}

El motor necesita una taxonomía de incidentes porque la prioridad no se interpreta igual en todos los dominios. Una llamada por incendio forestal requiere evaluar propagación, simultaneidad y amenaza a población o infraestructuras. Un accidente de tráfico exige valorar atrapamiento, heridos, accesibilidad y número de afectados. Una fuga química requiere identificar intoxicación, transporte de mercancías peligrosas y potencial de escalado. Por tanto, la categoría preliminar no es una etiqueta decorativa, sino una entrada que ayuda a activar las reglas correctas.

En v0.1.0, la taxonomía se deriva de dos fuentes: el catálogo de riesgos y categorías operativas recogido en la normativa, y los campos textuales del Registro 112 CyL. El campo estructurado de tipo de incidente resulta útil como orientación inicial, pero el sistema no debe depender exclusivamente de él. La descripción textual y el título del incidente son necesarios para detectar señales como \texttt{herido grave}, \texttt{atrapado}, \texttt{incendio}, \texttt{rescate}, \texttt{intoxicación} o \texttt{varias llamadas}.

La taxonomía cumple tres funciones. Primero, reduce ambigüedad al normalizar incidentes heterogéneos. Segundo, conecta cada incidente con variables operativas V01--V15. Tercero, mejora la explicabilidad, porque permite justificar la prioridad con factores coherentes con la naturaleza del suceso. Esta taxonomía se desarrolla en el Anexo A y se utiliza de forma transversal en los capítulos de diseño, datos y evaluación.

\section{Fuentes de datos y corpus normativo}

El prototipo combina dos tipos de fuentes. Por un lado, el Registro de Emergencias del 112 de Castilla y León aporta incidentes reales, texto operativo, fechas, localización y campos útiles para probar la ingesta, limpieza y extracción de señales \cite{Registro112CyL}. Por otro, el corpus normativo CyL alimenta la capa RAG y las herramientas MCP encargadas de recuperar citas legales, explicar reglas y sostener normativamente las recomendaciones.

La Tabla~\ref{tab:cap3_fuentes} resume el rol de las fuentes principales en el sistema.

\begin{table}[htbp]
\centering
\small
\begin{tabular}{|p{3.4cm}|p{4.1cm}|p{5.5cm}|}
\hline
\textbf{Fuente} & \textbf{Rol en el TFM} & \textbf{Uso principal} \\ \hline
Registro 112 CyL & Fuente empírica real & Ingesta, limpieza, texto libre, coordenadas, señales operativas y validación del pipeline \\ \hline
Ley 17/2015 y RD 524/2023 & Marco estatal & Conceptos de riesgo, emergencia, planificación, catálogo de riesgos y gradación de respuesta \\ \hline
Ley 4/2007 CyL y PLANCAL & Marco autonómico & Coordinación, inmediatez, guía de etiquetado P1--P4 y trazabilidad territorial \\ \hline
INFOCAL, INUNCYL, MPCyL & Riesgos específicos & Incendios forestales, inundaciones y mercancías peligrosas; en MPCyL solo consta el decreto de aprobación y queda pendiente el plan técnico completo \\ \hline
RGPD, LOPDGDD, Ley 11/2022 & Datos y comunicaciones & Minimización, localización, seguridad, logs y exclusión de variables sensibles o no necesarias \\ \hline
Reglamento UE 2024/1689 & IA responsable & Supervisión humana, transparencia, trazabilidad, gestión del riesgo y robustez \\ \hline
\end{tabular}
\caption{Fuentes principales y función dentro del sistema.}
\label{tab:cap3_fuentes}
\end{table}

El uso del Registro 112 CyL requiere una cautela esencial: el dataset no contiene una prioridad oficial P1--P4. Por ello, la etiqueta empleada en el proyecto es académica y se construye mediante supervisión débil a partir de varios anotadores, reglas y criterios derivados de la guía de etiquetado. Esta decisión evita presentar como verdad institucional una prioridad que no existe en la fuente original.

El corpus normativo, por su parte, no se utiliza para entrenar directamente el motor de prioridad. Su función principal es sostener la capa de explicación: buscar fragmentos relevantes, recuperar bases legales, citar normas y mostrar al operador por qué una regla o una prioridad tienen respaldo documental. Esta separación entre datos empíricos, reglas de priorización y fuentes normativas es clave para mantener la trazabilidad del sistema.

\section{Trazabilidad normativa de variables y reglas}

La Tabla~\ref{tab:cap3_trazabilidad} sintetiza la traducción del marco normativo a criterios de modelado. No debe leerse como una equivalencia jurídica exacta, sino como una matriz de trazabilidad: indica qué norma justifica cada grupo de variables o reglas dentro del prototipo. Esta estructura bidireccional facilita la auditoría del sistema por parte de expertos legales y técnicos, garantizando que cada elemento predictivo responda a una base regulatoria.

\begin{table}[htbp]
\centering
\scriptsize
\begin{tabular}{|p{3.1cm}|p{3.2cm}|p{3.5cm}|p{3.2cm}|}
\hline
\textbf{Norma / fuente} & \textbf{Criterio operativo} & \textbf{Variable o señal} & \textbf{Uso en el sistema} \\ \hline
Ley 17/2015 & Protección de la vida, daños personales, vulnerabilidad & V01, V02, V03, V05, \texttt{signal\_herido\_grave}, \texttt{signal\_atrapado} & Reglas duras P1 y aumento de prioridad \\ \hline
RD 524/2023 & Catálogo de riesgos y planificación & V04 & Taxonomía y normalización de tipo de incidente \\ \hline
PLEGEM & Coordinación ampliada y multirriesgo & V13, V14 & Señales de escalado y complejidad \\ \hline
Ley 4/2007 CyL & Coordinación, inmediatez y capacidad de respuesta & V07, V15 & Ajuste por emplazamiento crítico y accesibilidad \\ \hline
PLANCAL & Fases, situaciones y lógica territorial de respuesta & Escala P1--P4 académica & Guía de etiquetado y explicación \\ \hline
INFOCAL & Incendio forestal y propagación & V12, \texttt{signal\_incendio} & Reglas P2/P1 según riesgo de propagación \\ \hline
INUNCYL & Inundación y riesgo territorial & V09, V10, \texttt{signal\_meteo\_inundacion} & Anclaje v0.2.0; contexto RAG en v0.1.0 \\ \hline
MPCyL & Mercancías peligrosas y riesgo químico & \texttt{signal\_intoxicacion} & Regla de alta prioridad para fuga, intoxicación o cisterna; anclaje limitado por ausencia del plan técnico completo \\ \hline
Registro 112 CyL & Incidentes reales y texto operativo & Título, descripción, fecha, provincia, coordenadas & Fuente empírica; no contiene etiqueta P1--P4 \\ \hline
RGPD y LOPDGDD & Minimización y protección de datos & Política de exclusión & Logs, hashes, retención y no reidentificación \\ \hline
Reglamento UE 2024/1689 & Supervisión humana y trazabilidad & Arquitectura HITL, logs, explicaciones & Recomendación no vehemente y auditoría \\ \hline
\end{tabular}
\caption{Trazabilidad normativa entre criterios, variables y uso previsto en el prototipo.}
\label{tab:cap3_trazabilidad}
\end{table}

\section{Alcance v0.1.0 y elementos diferidos}

La excelencia del proyecto depende tanto de lo que se incluye como de lo que se decide diferir. En v0.1.0 se prioriza un núcleo evaluable: variables textuales y operativas centrales, RuleFit como motor interpretable, explicación con LLM local, RAG normativo, backend auditable y supervisión humana. Esta versión activa las variables V01--V07 y V12--V15, junto con las señales textuales necesarias para detectar riesgo vital, heridos, atrapamiento, intoxicación, incendio, rescate, accidente de tráfico, varias llamadas y riesgo meteorológico textual.

Quedan diferidas a v0.2.0 las variables que requieren integraciones contextuales más costosas: avisos AEMET, condición meteorológica adversa, zonas inundables SNCZI y proximidad a instalaciones Seveso. Estas variables mantienen su anclaje normativo en el capítulo y en el modelo de datos, pero no forman parte del núcleo entrenado y evaluado en v0.1.0. Esta decisión es metodológicamente prudente: evita introducir fuentes geoespaciales o históricas que podrían debilitar la reproducibilidad del TFM, sin renunciar a una línea clara de evolución futura.

La exclusión temporal de AEMET, SNCZI y Seveso es también una decisión de diseño orientada a preservar el aislamiento del entrenamiento local del modelo base. Incorporar dichas fuentes exigiría sincronizar servicios externos, resolver coberturas territoriales, transformar coordenadas, versionar capas geográficas y documentar criterios de vigencia temporal. En una primera versión académica, priorizar variables textuales y operativas disponibles en el Registro 112 CyL permite comparar baseline experto, RuleFit-lite e \texttt{imodels} bajo condiciones reproducibles. El principal riesgo de esta decisión es que la priorización automática pueda ser menos sensible ante inundaciones, meteorología adversa o riesgos tecnológicos específicos; por ello, estas integraciones se declaran formalmente como extensión v0.2.0.

\section{Conclusiones del capítulo}

El marco operativo y normativo analizado permite fijar cinco conclusiones clave para el desarrollo del resto de este Trabajo Fin de Máster. Estas deducciones metodológicas sintetizan la relación entre la práctica operacional y las limitaciones predictivas del sistema. Las lecciones extraídas se detallan en los siguientes puntos:

En primer lugar, el sistema debe entenderse como apoyo a la decisión, no como sustituto del operador. Esta conclusión deriva tanto de la naturaleza de la protección civil como del marco europeo de inteligencia artificial. De este modo, la responsabilidad última sobre las acciones operativas permanece siempre bajo la supervisión humana directa.

En segundo lugar, la prioridad operativa no puede reducirse a una única señal. La gravedad actual, la vulnerabilidad, la exposición, la propagación, la simultaneidad, la accesibilidad y la incertidumbre informativa deben combinarse mediante contratos y reglas trazables. Esta integración multifactorial asegura que la clasificación del incidente sea representativa de su complejidad real.

En tercer lugar, la escala P1--P4 es académica y funcional. Se inspira en la lógica de gradación de la normativa, especialmente PLANCAL, pero no equivale a situaciones operativas oficiales ni pretende reemplazarlas. Se trata de una herramienta de simulación científica adecuada para evaluar el comportamiento de los modelos matemáticos propuestos.

En cuarto lugar, el Registro 112 CyL aporta realismo empírico, pero no una verdad supervisada de prioridad. Por eso el proyecto necesita supervisión débil, validación interna y comparación explícita entre baseline experto, RuleFit-lite e \texttt{imodels}. Esta aproximación honesta frente al dataset original garantiza el rigor científico en la posterior medición de resultados.

En quinto lugar, la trazabilidad normativa no es un adorno documental: es parte de la arquitectura. Cada recomendación debe poder conectarse con variables, reglas, citas legales, versión de modelo y decisión humana final. Esta exigencia enlaza directamente con los capítulos siguientes, donde se desarrollan la metodología, la arquitectura de tres capas, el diseño de datos y la evaluación del prototipo.
